\section{Square prefixes and the arithmetic decomposition}
\label{sec:arithmetic}
%======================================================================

The square-root cutoff is the arithmetic hinge of the paper.  Before introducing geometry, we isolate the exact decomposition that the geometry is meant to explain.

Define the complete square block
\begin{equation}
 B_n:=\{m\in\N:n^2\le m<(n+1)^2\},
 \label{eq:block}
\end{equation}
the square-root cutoff
\begin{equation}
 R_n:=n+1,
 \label{eq:Rn}
\end{equation}
and the complete-prefix endpoint
\begin{equation}
 X_n:=R_n^2-1=(n+1)^2-1.
 \label{eq:Xn}
\end{equation}
The square-prefix Mertens value is
\begin{equation}
 S_n:=M(X_n)=\sum_{m\le X_n}\mu(m).
 \label{eq:Sn}
\end{equation}
Because $\mu(R_n^2)=0$ for $R_n>1$, one may equally write $S_n=M(R_n^2)$.

Write $\Pplus(m)$ for the largest prime factor of $m$, with $\Pplus(1)=1$.

\begin{definition}[Smooth and transport terms]
\label{def:AT}
Define
\begin{align}
 A_n
 &:={\sum}_{\substack{1\le m\le X_n\\\Pplus(m)\le R_n}}\mu(m),
 \label{eq:An}\\
 T_n
 &:={\sum}_{1\le c<R_n}\mu(c)
 \left[
 \pi\!\left(\left\lfloor\frac{X_n}{c}\right\rfloor\right)-\pi(R_n)
 \right].
 \label{eq:Tn}
\end{align}
\end{definition}

\begin{lemma}[One-large-prime closure]
\label{lem:one-large}
Let $m\le R^2-1$ be squarefree.  If $\Pplus(m)>R$, then there is a unique factorization
\begin{equation}
 m=cq,
 \qquad
 1\le c<R<q,
 \qquad
 q\text{ prime},
 \label{eq:one-large}
\end{equation}
where $q=\Pplus(m)$.  Moreover,
\begin{equation}
 \mu(m)=-\mu(c).
 \label{eq:mobius-peel}
\end{equation}
\end{lemma}

\begin{proof}
Put $q=\Pplus(m)$ and $c=m/q$.  If $c\ge R$, then $m=cq>R^2$, contrary to $m\le R^2-1$.  Since $m$ is squarefree, $q\nmid c$ and $\mu(cq)=-\mu(c)$.  Uniqueness follows from the largest-prime choice.
\end{proof}
\leanxref{Lean support}{canonical largest-prime and cofactor arithmetic}{https://github.com/OVVO-Financial/RH_Square_Block/blob/main/lean/RHLean/Analysis/CanonicalHighSectorCore.lean}

\begin{proof}[Proof of \cref{thm:intro-decomposition}]
Partition the squarefree $m\le X_n$ according to whether $\Pplus(m)\le R_n$ or $\Pplus(m)>R_n$.  The first population contributes $A_n$.  By \cref{lem:one-large}, the second population is uniquely $m=cq$ with $c<R_n<q$, and each source has weight $-\mu(c)$.  The total second contribution is therefore $-T_n$.
\end{proof}

\begin{corollary}[Elementary smooth--transport RH inequality]
\label{cor:elementary-AT-RH}
The Riemann Hypothesis is equivalent to either of the following square-prefix estimates:
\begin{align}
 |A_N-T_N|^2&\ll_\eps N^{2+\eps},
 \label{eq:elementary-AT-squared}\\
 |A_N-T_N|&\ll_\eps N^{1+\eps}.
 \label{eq:elementary-AT-linear}
\end{align}
Here and below the two exponent conventions are equivalent after renaming $\eps$.
\end{corollary}

\begin{proof}
By \cref{thm:intro-decomposition}, $A_N-T_N=M((N+1)^2-1)$.  The classical Mertens criterion gives the estimates under RH.  Conversely, every $x$ lies within one square block of an endpoint $(N+1)^2-1$, and $|\mu(m)|\le1$, so the endpoint estimate interpolates to $M(x)\ll_\eps x^{1/2+\eps}$.
\end{proof}

\begin{remark}[Why this is the most elementary total formulation]
The inequality \cref{eq:elementary-AT-linear} contains no shell, window, covariance, probability space, or operator.  It asks only that the two exact cumulative arithmetic terms in $S_N=A_N-T_N$ agree to square-root scale.  The high-sector formulation introduced later removes the part of this agreement already forced by elementary geometry.
\end{remark}

\subsection{The largest-prime cofactor identity}
\label{subsec:central-identity}

For $c\ge1$, define the prime-tail counting function
\begin{equation}
 \Pi_c(x)
 :=
 \#\left\{
 q\text{ prime}:
 \Pplus(c)<q\le\left\lfloor\frac{x}{c}\right\rfloor
 \right\}.
 \label{eq:Pi-c}
\end{equation}
Equivalently,
\begin{equation}
 \Pi_c(x)
 =
 \left(
 \pi\!\left(\left\lfloor\frac{x}{c}\right\rfloor\right)
 -\pi(\Pplus(c))
 \right)_{+},
 \label{eq:Pi-c-positive-part}
\end{equation}
where $(u)_+:=\max\{u,0\}$.  The positive part is essential when
$x/c<\Pplus(c)$.

\begin{theorem}[Exact largest-prime cofactor expansion]
\label{thm:central-M-identity}
For every real $x\ge1$,
\begin{equation}
 \boxed{
 M(x)
 =
 1-\sum_{c\ge1}\mu(c)\Pi_c(x).
 }
 \label{eq:central-M-identity-compact}
\end{equation}
The sum is finite, and separating the prime curve $c=1$ gives
\begin{equation}
 \boxed{
 M(x)
 =
 1-\pi(x)
 -\sum_{2\le c\le x/2}
 \mu(c)
 \left(
 \pi\!\left(\left\lfloor\frac{x}{c}\right\rfloor\right)
 -\pi(\Pplus(c))
 \right)_{+}.
 }
 \label{eq:central-M-identity}
\end{equation}
\end{theorem}

\begin{proof}
Every squarefree $m>1$ has a unique largest prime factor
$q=\Pplus(m)$.  Writing $c=m/q$, one has
\[
 q>\Pplus(c),
 \qquad
 \mu(m)=\mu(cq)=-\mu(c).
\]
Conversely, every pair consisting of a squarefree $c$ and a prime
$q>\Pplus(c)$ produces the squarefree integer $m=cq$ whose largest
prime factor is $q$.  The condition $m\le x$ is exactly $q\le x/c$.
Summing $-\mu(c)$ over these unique pairs proves
\cref{eq:central-M-identity-compact}.  The $c=1$ term is
$-\pi(x)$, because $\Pplus(1)=1$.  No term with $c>x/2$ can occur.
\end{proof}

\begin{theorem}[Extreme largest-prime support and the pure-prime fibre]
\label{thm:extreme-largest-prime}
Let $x\ge2$ and let $1<m\le x$ be an integer.  Put
\[
 q=\Pplus(m),\qquad c=\frac{m}{q}.
\]
If
\begin{equation}
 x<2q,
 \label{eq:extreme-q-condition}
\end{equation}
then
\begin{equation}
 \boxed{c=1,\qquad m=q,}
 \label{eq:extreme-prime-fibre}
\end{equation}
so $m$ is prime.  Equivalently, every composite integer $m\le x$ satisfies
\begin{equation}
 \boxed{\Pplus(m)\le\left\lfloor\frac{x}{2}\right\rfloor.}
 \label{eq:composite-largest-prime-half}
\end{equation}
Moreover, for every genuine canonical product $m=cq\le x$ with $c\ge2$ one also has
\begin{equation}
 \boxed{q\le\left\lfloor\frac{x}{2}\right\rfloor,\qquad
 c\le\left\lfloor\frac{x}{2}\right\rfloor.}
 \label{eq:both-canonical-factors-half}
\end{equation}
Thus canonical sources below $x$ split sharply into the pure prime fibre $c=1$ and genuine product fibres $c\ge2$; the extreme-largest-prime region contains no composites.
\end{theorem}

\begin{proof}
Since $q\mid m$, write $m=cq$ with $c\ge1$.  If $c\ge2$, then
\[
 2q\le cq=m\le x,
\]
contradicting $x<2q$.  Hence $c=1$, so $m=q$ is prime.  The contrapositive gives \cref{eq:composite-largest-prime-half}.  For a genuine product $c\ge2$, the inequality $2q\le cq\le x$ gives the bound on $q$; since every canonical prime satisfies $q\ge2$, the inequality $2c\le cq\le x$ gives the bound on $c$.
\end{proof}
\leanxref{Lean proof}{canonicalCofactor_eq_one_of_endpoint_lt_two_mul_largestPrime; canonicalLargestPrimeFactor_le_half_of_ne_largestPrime}{https://github.com/OVVO-Financial/RH_Square_Block/blob/main/lean/RHLean/Analysis/CanonicalExtremePrimeSupport.lean}

\begin{remark}[Two different half-support facts]
The bounds in \cref{thm:extreme-largest-prime} should not be conflated.  The already-used cofactor support $c\le\lfloor x/2\rfloor$ follows from the trivial fact that the canonical prime is at least $2$.  The new extreme-prime statement says that a \emph{composite} source also has $q\le\lfloor x/2\rfloor$.  The only sources whose largest prime exceeds that threshold are the primes themselves on the $c=1$ fibre.
\end{remark}

\begin{remark}[Central analytic engine]
\Cref{eq:central-M-identity} is the exact arithmetic form of the
cofactor-parabola decomposition.  It shows that the Mertens function is
not obtained by proving that the prime curve is small.  The prime curve
$-\pi(x)$ is cancelled by the M\"obius-weighted collection of all
curves $c\ge2$.
\end{remark}

